\documentclass[11pt]{article}
\usepackage[margin=1in]{geometry}\usepackage[T1]{fontenc}\usepackage[utf8]{inputenc}
\usepackage{lmodern,amsmath,amssymb,amsthm,mathtools,booktabs,microtype}
\usepackage[colorlinks=true,linkcolor=blue,citecolor=blue,urlcolor=blue]{hyperref}
\hypersetup{pdftitle={A path-variation prior in Gaussian Bayesian ReLU networks},pdfauthor={ChatGPT-6 Astra; Claude Opus 5.5; Kaarel Hänni; Dmitry Vaintrob}}
\newtheorem{theorem}{Theorem}[section]\newtheorem{lemma}[theorem]{Lemma}
\newtheorem{corollary}[theorem]{Corollary}\newtheorem{proposition}[theorem]{Proposition}
\theoremstyle{definition}\newtheorem{definition}[theorem]{Definition}
\theoremstyle{remark}\newtheorem{remark}[theorem]{Remark}
\newcommand{\R}{\mathbb R}\newcommand{\E}{\mathbb E}\newcommand{\Prob}{\mathbb P}
\newcommand{\Rad}{\mathfrak R}\newcommand{\KL}{\operatorname{KL}}
\newcommand{\op}{\mathrm{op}}\newcommand{\VarPath}{\mathcal V}
\title{A path-variation prior\\in Gaussian Bayesian ReLU networks}
\author{ChatGPT-6 Astra \and Claude Opus 5.5 \and Kaarel H\"anni \and Dmitry Vaintrob\thanks{Authors are listed in alphabetical order.}}
\date{17 November 2026}
\begin{document}\maketitle
\begin{abstract}
We construct a finite Gaussian Bayesian ReLU learner whose induced function prior is quantitatively controlled by minimal path variation. Each learner has one fixed dense architecture. Internal matrices are conditioned to have absolute spectral norm at most one; positive homogeneity permits every ReLU function to satisfy this condition after neuron rescaling, without changing width or depth. A path-flow balancing identity makes the minimum normalized quadratic energy converge exactly to twice minimal path variation as learner width grows. Precision-normalized negative log probabilities of uniform function neighborhoods obey corresponding two-sided bounds. Labels may follow any bounded IID law. With a Gaussian likelihood calibrated from observed free energy, the ordinary raw posterior mean satisfies an agnostic oracle inequality. A teacher with path variation $P$, hidden depth $H$, and population squared loss $R_*$ needs $\widetilde O(H(1+P)^2(R_*/\varepsilon^2+1/\varepsilon))$ samples for excess at most $\varepsilon$. Inputs lie in a Euclidean unit ball; the bound is independent of input dimension at every learner width above the teacher width. A stable width-$C$ teacher with bounded endpoint norms has $P=O(C)$, giving quadratic width and linear depth. General small networks can have much larger path variation. The prior uses explicitly different endpoint and internal Gaussian scales and does not mix over network sizes.
\end{abstract}

\section{What the prior should prefer}
A small neural network may approximate a regression function well even when no small network fits it exactly. A useful Bayesian guarantee should compare test loss with such approximations. This paper supplies an explicit dense Gaussian learner with an approximation guarantee controlled by a standard network quantity: path variation.

The conclusion is also a statement about the induced prior on functions. At a fixed approximation scale, negative log prior probability is bounded above and below by minimal path variation, after normalization by the Gaussian precision and up to logarithmic terms. An exact variational identity describes the associated quadratic energy. The learning guarantee is a consequence of this structure and a capacity estimate, rather than a definition of the phrase ``path-variation prior.''

We use a norm condition on internal parameters. It is stronger than ordinary spectral nonexpansiveness because it does not rely on cancellation of signs. For ReLU, however, the condition removes no representable functions at a fixed width and depth: a positive rescaling of individual neurons enforces it. This reparametrization can change coefficient sizes and prior probabilities, which is precisely relevant to the prior-ordering result.

\section{Architecture, input convention, and path variation}
Let $\mathcal X$ be a compact subset of the Euclidean unit ball in $\R^d$, $d\ge1$. There are $H\ge1$ hidden ReLU layers, a scalar linear readout, and common hidden width $N$:
\[
 h^1(x)=\phi(W^1x),\qquad h^\ell(x)=\phi(W^\ell h^{\ell-1}(x)),\quad
 z_w(x)=v^{\mathsf T}h^H(x),\qquad\phi(t)=\max\{t,0\}.
\]
The first matrix is $N$ by $d$ and the internal matrices are $N$ by $N$. A constant input coordinate may be included in $x$. For ordinary first-layer biases on a unit-ball input $x_0$, use $x=(1,x_0)/\sqrt2$ and count the extra coordinate in $d$. There are no internal biases. Total affine depth is $L=H+1$.

For a network of any finite width, let $a^1_i=\|W^1_{i,:}\|_2$. Its mixed path variation is
\begin{equation}
 P(w)=|v|^{\mathsf T}|W^H|\cdots|W^2|a^1.
 \label{eq:pathnorm}
\end{equation}
For $H=1$, the intervening product is empty. Absolute values on matrices and vectors are entrywise. The Euclidean norm at the first layer matches the input unit ball. In particular $|z_w(x)|\le P(w)$ on $\mathcal X$.

For a function $f$ realizable at depth $H$, define $\VarPath_H(f)$ as the infimum of $P(w)$ over all finite-width depth-$H$ implementations of $f$ on $\mathcal X$. For a width bound $N$ and accuracy $\delta\ge0$, define
\begin{equation}
 \VarPath_{H,N,\delta}(f)=
 \inf\{P(w):\operatorname{width}(w)\le N,\ \|z_w-f\|_{\infty,\mathcal X}\le\delta\}.
 \label{eq:approx-variation}
\end{equation}
The infimum of an empty set is infinity. These are implementation complexity measures, not probabilities.

\section{The Gaussian prior}
The internal support condition is
\begin{equation}
 \bigl\||W^\ell|\bigr\|_{\op}\le1\qquad(2\le\ell\le H).
 \label{eq:support}
\end{equation}
For any $\gamma\ge1$, put $\omega=N(d+1)+(H-1)N^{\gamma+1}$. At $\gamma=1$, this is exactly the number of Gaussian parameter coordinates. Before conditioning on~\eqref{eq:support}, all Gaussian coordinates are independent and centered, with
\begin{equation}
 \operatorname{Var}(W^1_{ji})=\frac N\omega,\qquad
 \operatorname{Var}(W^\ell_{ji})=N^{-\gamma}\ (\ell\ge2),\qquad
 \operatorname{Var}(v_j)=\frac1{\omega N}.
 \label{eq:prior}
\end{equation}
The input convention is why no factor $1/d$ appears in the first variance. Denote the conditioned prior by $\pi_{N,\gamma}$.

Define endpoint energy and total normalized energy by
\begin{align}
 u_N&=\frac1N\|W^1\|_F^2,\qquad q_N=N\|v\|_2^2,\qquad Q_N=u_N+q_N,\notag\\
 E_N(w)&=Q_N(w)+\frac{N^\gamma}{\omega}\sum_{\ell=2}^H\|W^\ell\|_F^2.
 \label{eq:energy}
\end{align}
The prior density on its support is proportional to $e^{-\omega E_N(w)/2}$. The endpoint energy is independent of this normalization; the internal coefficient is at most $1/((H-1)N)$ when $H>1$.

To compare with a teacher of width $C$, the only width requirement is $N\ge C$; there is no additional input-dimension threshold. There is no distribution over learner widths and no limiting operation in those results.

\section{Path-flow balancing and the exact energy characterization}
\begin{lemma}[Path-flow balancing]
\label{lem:flow}
Every finite ReLU network with path variation $P>0$ has an equivalent implementation at the same width and depth such that every internal absolute spectral norm is at most one and
\begin{equation}
 \|W^1\|_F^2=P,\qquad \|v\|_2^2=P.
 \label{eq:balanced-flow}
\end{equation}
The path variation is unchanged. If $P=0$, the function is zero.
\end{lemma}
\begin{proof}
Write $A^\ell=|W^\ell|$. Starting from $a^1_i=\|W^1_{i,:}\|_2$, set $a^\ell=A^\ell a^{\ell-1}$. Starting from $b^H=|v|$, set $b^{\ell-1}=(A^\ell)^{\mathsf T}b^\ell$. At every hidden layer,
\[
 \sum_j a^\ell_jb^\ell_j=P.
\]
Neurons with zero forward or backward mass lie on no nonzero input-to-output path and can be disconnected without changing the function. On the remaining neurons define
\[
 s^\ell_j=\sqrt{b^\ell_j/a^\ell_j},\qquad
 r^\ell_j=\sqrt{a^\ell_jb^\ell_j}.
\]
Multiply hidden activation $j$ in layer $\ell$ by $s^\ell_j$ by rescaling its incoming weights and inversely rescaling its outgoing weights. ReLU positive homogeneity preserves the exact function. The transformed internal absolute matrix is
$\widetilde A^\ell=\operatorname{diag}(s^\ell)A^\ell\operatorname{diag}(s^{\ell-1})^{-1}$, and it satisfies
\[
 \widetilde A^\ell r^{\ell-1}=r^\ell,\qquad
 (\widetilde A^\ell)^{\mathsf T}r^\ell=r^{\ell-1}.
\]
Weighted Cauchy--Schwarz gives, for any vector $x$ on the active coordinates,
\begin{align*}
 \sum_j\left(\sum_i\widetilde A^\ell_{ji}x_i\right)^2
 &\le\sum_j r^\ell_j\sum_i\widetilde A^\ell_{ji}\frac{x_i^2}{r^{\ell-1}_i}\\
 &=\sum_i\frac{x_i^2}{r^{\ell-1}_i}
                ((\widetilde A^\ell)^{\mathsf T}r^\ell)_i
 =\sum_i x_i^2.
\end{align*}
Thus its operator norm is at most one. The transformed first-row Euclidean norm is $s^1_ja^1_j=r^1_j$, and the transformed readout magnitude is $b^H_j/s^H_j=r^H_j$. Their squared sums are both $P$, proving~\eqref{eq:balanced-flow}. Positive neuron rescaling cancels on every path, so path variation is preserved. If $P=0$, every absolute path product is zero, and $|z_w|\le P$ gives the zero function.
\end{proof}

\begin{lemma}[Cloning without a width divisibility condition]
\label{lem:clone}
A teacher of width $C$ and path variation $P$ has an exact supported implementation at every $N\ge C$ with
\begin{equation}
 Q_N=2P,\qquad
 \sum_{\ell=2}^H\|W^\ell\|_F^2\le(H-1)C.
 \label{eq:clone}
\end{equation}
For the zero function, the all-zero implementation suffices.
\end{lemma}
\begin{proof}
First apply Lemma~\ref{lem:flow}. Put $k=\lfloor N/C\rfloor$, clone every teacher neuron exactly $k$ times, and leave the other neurons disconnected with zero weights. The internal blocks are $W^\ell\otimes(J_k/k)$. Both their absolute spectral norms and their squared Frobenius norms are unchanged. Since a $C$ by $C$ matrix of operator norm at most one has squared Frobenius norm at most $C$, the internal bound follows.

The endpoint energies before the final balancing are $u_N=(k/N)P$ and $q_N=(N/k)P$. Multiply the first layer by $\sqrt{N/k}$ and divide the readout by the same factor. Positive homogeneity again preserves the function, and now $u_N=q_N=P$. This proves~\eqref{eq:clone} exactly, including widths not divisible by $C$.
\end{proof}

\begin{theorem}[An exact variational path-prior characterization]
\label{thm:energy-characterization}
For every supported network,
\begin{equation}
 E_N(w)\ge Q_N(w)\ge2P(w).
 \label{eq:energy-lower}
\end{equation}
Conversely, a teacher of width $C$ and path variation $P$ has, at each $N\ge C$, an exact supported implementation with
\begin{equation}
 E_N\le2P+\frac{N^\gamma(H-1)C}{\omega}\le2P+\frac CN.
 \label{eq:energy-upper}
\end{equation}
Consequently, for every finite-network-realizable function $f$ at depth $H$,
\begin{equation}
 \lim_{N\to\infty}\inf_{\substack{w\text{ supported at width }N\\z_w=f\text{ on }\mathcal X}}E_N(w)
 =2\VarPath_H(f).
 \label{eq:energy-limit}
\end{equation}
\end{theorem}
\begin{proof}
Under the support condition,
\[
 P(w)\le\|v\|_2\left(\sum_i\|W^1_{i,:}\|_2^2\right)^{1/2}
          =\sqrt{u_Nq_N}\le Q_N/2.
\]
This proves~\eqref{eq:energy-lower}; Lemma~\ref{lem:clone} proves~\eqref{eq:energy-upper}. The lower bound in~\eqref{eq:energy-limit} holds at every width. For any $\epsilon>0$, choose a finite-width implementation of $f$ with path variation at most $\VarPath_H(f)+\epsilon$. Its width is fixed while $N$ grows, so~\eqref{eq:energy-upper} gives upper limit at most $2\VarPath_H(f)+2\epsilon$. Let $\epsilon$ decrease to zero.
\end{proof}

\section{An exact conditional-Gaussian comparison}
The set $\{W:\||W|\|_{\op}\le1\}$ is convex, centrally symmetric, and has positive Gaussian measure. Convexity follows from entrywise absolute-value convexity and monotonicity of the operator norm on nonnegative matrices.

We will use the following identity. If $X$ has density proportional to $\mathbf1_K(x)e^{-\alpha\|x\|^2/2}$ on a convex centrally symmetric set $K\subset\R^p$, $a\in K$, and $Y=(1-\eta)a+\eta X$ with $0<\eta<1$, then
\begin{align}
 \KL(\mathcal L(Y)\|\mathcal L(X))
 &=p\log(1/\eta)+\frac\alpha2\left[(1-\eta)^2\|a\|^2-(1-\eta^2)\E\|X\|^2\right]\notag\\
 &\le p\log(1/\eta)+\frac\alpha2\|a\|^2.
 \label{eq:transport}
\end{align}
Indeed, $Y\in K$, the affine Jacobian is $\eta^p$, and symmetry gives $\E X=0$. The conditional-Gaussian normalizer cancels exactly in the density ratio. For an unrestricted Gaussian, $Y=a+\eta X$ obeys the same upper bound.

\begin{lemma}[Uniform approximation with controlled relative entropy]
\label{lem:comparison}
Let $g$ have width $C$, depth $H$, and mixed path variation $P$. At every $N\ge C$, every $\gamma\ge1$, and every $0<\eta<1$, there is a law $\rho_{g,\eta}$ relative to $\pi_{N,\gamma}$ such that
\begin{equation}
 \frac2\omega\KL(\rho_{g,\eta}\|\pi_{N,\gamma})
 \le2P+1+2\log(1/\eta),
 \label{eq:comparison-entropy}
\end{equation}
and
\begin{equation}
 \left(\int\|z_w-g\|_{\infty,\mathcal X}^2d\rho_{g,\eta}\right)^{1/2}
 \le2(H+1)(P+1)\eta.
 \label{eq:comparison-uniform}
\end{equation}
\end{lemma}
\begin{proof}
Take the exact implementation from Lemma~\ref{lem:clone}, with endpoint energies $u_N=q_N=P$. Independently for every internal layer, draw $X$ from its actual conditioned Gaussian prior and set $W^\ell=(1-\eta)W^{\ell,*}+\eta X$. For the first layer and readout use Gaussians with the cloned means and $\eta^2$ times their original covariances.

The dimension cost in~\eqref{eq:transport}, normalized by $2/\omega$, is at most $2\log(1/\eta)$ because
\[
 \frac{Nd+(H-1)N^2+N}{N(d+1)+(H-1)N^{\gamma+1}}\le1.
\]
The endpoint mean cost is exactly $2P$. The internal mean cost is at most $N^\gamma(H-1)C/\omega\le C/N\le1$. This proves~\eqref{eq:comparison-entropy}.

Let $e_\ell$ be the $L^2(\rho)$ norm of the maximum over $x\in\mathcal X$ of $\|h^\ell(x)-h^{\ell,*}(x)\|_2/\sqrt N$. The first Gaussian perturbation has
\[
 \frac1N\E\|W^1-W^{1,*}\|_F^2
 =\eta^2\frac{Nd}{\omega}\le\eta^2,
\]
so $e_1\le\eta$. Every internal random matrix has operator norm at most one and differs from its teacher matrix by operator norm at most $2\eta$. Teacher hidden states satisfy $\sup_x\|h^{\ell,*}(x)\|_2/\sqrt N\le\sqrt P$. Consequently
\[
 e_H\le\eta+2(H-1)\eta\sqrt P.
\]
The cloned readout contributes at most $\sqrt P\,e_H$ to the output error.

The readout noise is independent of the hidden layers and has expected squared Euclidean norm $\eta^2/\omega$. Since all internal layers are Euclidean nonexpansive,
\begin{align*}
 \E\sup_x| (v-v^*)^{\mathsf T}h^H(x)|^2
 &\le\frac{\eta^2}{\omega}\E\|W^1\|_F^2\\
 &=\eta^2\frac N\omega\left(P+\eta^2\frac{Nd}\omega\right)
 \le\eta^2(P+1).
\end{align*}
Minkowski gives total error at most
$\eta\{\sqrt P+2(H-1)P+\sqrt{P+1}\}\le2(H+1)(P+1)\eta$.
This is a uniform-function estimate, not merely a pointwise-in-input estimate.
\end{proof}

For learning, put
\begin{equation}
 \eta_n=\frac1{16n(H+1)(P+1)^2},\qquad
 T_g=4+2P+2\log(1/\eta_n).
 \label{eq:teacher-budget}
\end{equation}
Then $|g|\le P$, and~\eqref{eq:comparison-uniform} makes the expected squared-loss increase at every observation at most $1/(2n)$. Thus the Gibbs free energy for observation variance $n\lambda/\omega$ satisfies
\begin{equation}
 F_S(\lambda)\le\widehat R_S(g)+\frac1{2n}+\lambda T_g
 \qquad(\lambda>0).
 \label{eq:free-energy-comparison}
\end{equation}
The slack from $2P+1$ to $4+2P$ is convenient for the population-risk constants below.

\section{The induced prior on function neighborhoods}
For a bounded function $f$ on $\mathcal X$ and $0<\delta<1$, write
\[
 B_\delta(f)=\{w:\|z_w-f\|_{\infty,\mathcal X}\le\delta\}.
\]
These events are measurable because the network map is continuous in parameters in the uniform norm on the compact input domain.

\begin{theorem}[Two-sided finite function-prior bounds]
\label{thm:prior}
For every $N\ge1$ and $\gamma\ge1$,
\begin{equation}
 -\frac1\omega\log\pi_{N,\gamma}(B_\delta(f))
 \ge\frac12\VarPath_{H,N,\delta}(f)-\frac{\log2}{2}.
 \label{eq:prior-lower}
\end{equation}
If $p=\VarPath_{H,N,\delta/2}(f)<\infty$, then
\begin{equation}
 -\frac1\omega\log\pi_{N,\gamma}(B_\delta(f))
 \le2p+1+2\log\frac{8(H+1)(p+1)}\delta+2\log2.
 \label{eq:prior-upper}
\end{equation}
In particular the precision-normalized negative log prior is controlled by minimal approximate path variation, with only the displayed logarithmic accuracy and depth terms. No learner-width or input-dimension factor is suppressed.
\end{theorem}
\begin{proof}
Support conditioning leaves the first layer and readout Gaussian and independent, so
\begin{equation}
 \omega Q_N\sim\chi^2_{N(d+1)},\qquad
 \frac{N(d+1)}\omega\le1.
 \label{eq:chi-square}
\end{equation}
Any network in $B_\delta(f)$ has $P(w)\ge\VarPath_{H,N,\delta}(f)$ and $Q_N\ge2P(w)$ by~\eqref{eq:energy-lower}. Markov's inequality with
$\E e^{\omega Q_N/4}=2^{N(d+1)/2}$ yields
\[
 \pi(B_\delta(f))\le
 \exp\left[-\frac\omega2\VarPath_{H,N,\delta}(f)
                      +\frac{N(d+1)\log2}{2}\right],
\]
which proves~\eqref{eq:prior-lower}. An infinite infimum means the event is empty.

For the upper bound, choose any width-at-most-$N$ network $g$ within $\delta/2$ of $f$, with path variation $P$. Apply Lemma~\ref{lem:comparison} at
$\eta=\delta/[8(H+1)(P+1)]$. Its uniform-error $L^2$ norm is at most $\delta/4$, so its comparison law gives probability at least $3/4$ to $B_{\delta/2}(g)\subset B_\delta(f)$. The binary relative-entropy inequality implies, for any event $A$ with $\rho(A)\ge1/2$,
\[
 \log\frac1{\pi(A)}\le2\KL(\rho\|\pi)+2\log2.
\]
Use~\eqref{eq:comparison-entropy}, divide by $\omega\ge1$, and let $P$ decrease to the infimum $p$. This proves~\eqref{eq:prior-upper}.
\end{proof}

The different radii in the two bounds are intentional: approximation at half the requested accuracy leaves room for an open neighborhood of the comparison implementation. The theorem is about function neighborhoods, not positive mass on one exact continuous-domain function. It gives bounds on prior preference; it does not assert an exact monotone ordering between every pair of functions.

\section{A dimension-free capacity bound}
Let $\mathcal F_t$ be the union over all positive widths $N$ of supported networks with $Q_N\le t$, at fixed hidden depth. Put
\[
 K_H=1+\sqrt{2H\log2}.
\]

\begin{lemma}[Capacity and amplitude]
\label{lem:capacity}
On every input sample in $\mathcal X$,
\[
 \widehat\Rad_n(\mathcal F_t)\le\frac{tK_H}{\sqrt n},\qquad
 |z_w(x)|\le P(w)\le Q_N(w)/2.
\]
The same capacity bound holds for worst-case samples, with no dimension dependence.
\end{lemma}
\begin{proof}
The amplitude statement follows from~\eqref{eq:energy-lower}. To bound capacity, normalize every hidden feature by its forward path mass: divide $h_j^\ell$ by $a_j^\ell$, discarding zero-mass coordinates. ReLU homogeneity makes every first-row Euclidean norm at most one, every internal absolute row sum at most one, and the transformed readout's absolute coefficient sum equal to $P(w)\le t/2$.

Let $\mathcal A_0$ be the sample vectors of all Euclidean-unit-norm linear functions. Recursively let $\mathcal A_\ell$ be the closed absolute convex hull of the coordinatewise ReLU images of $\mathcal A_{\ell-1}$. Every normalized hidden feature belongs to its corresponding class. For $M_\ell=\sup_{f\in\mathcal A_\ell}\sum_i\epsilon_if_i$ and $s>0$, exponential scalar contraction gives
\[
 \E e^{sM_\ell}\le2\E e^{sM_{\ell-1}}.
\]
Indeed, absolute convexification gives a maximum of positive and negative suprema. Its exponential is at most their summed exponentials, which have equal expectation. Contraction for either supremum has constant one: conditional on other signs, Lipschitzness controls the sum of the two coordinate suprema, and $|\phi(t)|\le|t|$ controls their maximum; these imply the inequality for the sum of exponentials.

At the initial layer, $M_0=\|\sum_i\epsilon_iX_i\|_2$. Jensen gives $\E M_0\le\sqrt n$. Changing sign $i$ changes $M_0$ by at most $2\|X_i\|_2\le2$, so the bounded-difference exponential inequality gives
\[
 \E e^{sM_0}\le e^{s\sqrt n+s^2n/2}.
\]
This last inequality also follows directly by applying Hoeffding's lemma to the Doob martingale of $M_0$. Iterating the contraction, taking logarithms, and optimizing $s$ yields $\E M_H\le\sqrt nK_H$. Multiply by the readout absolute sum at most $t/2$ and divide by $n$; the displayed capacity bound leaves a harmless factor two of slack. Separability follows from countable dense parameter subclasses at each width and a countable union over widths.
\end{proof}

\section{A data-dependent likelihood with a raw-risk certificate}
Let $R(f)=\E(f(X)-Y)^2$ and $\widehat R_S(f)=n^{-1}\sum_i(f(X_i)-Y_i)^2$. With Gaussian observation variance $n\lambda/\omega$, define
\[
 \Pi_{\lambda,S}(dw)=Z^{-1}e^{-\omega\widehat R_S(z_w)/(2\lambda)}\pi_{N,\gamma}(dw),\qquad
 F_S(\lambda)=-\frac{2\lambda}{\omega}\log Z.
\]
The variational identity is $F_S(\lambda)=\inf_\rho\{\int\widehat R_Sd\rho+2\lambda\KL(\rho\|\pi)/\omega\}$. The centered independent readout cancels cross terms and
\[
 \E_\pi z_w(x)^2\le\frac{\E\|W^1x\|_2^2}{\omega N}
                 \le\frac N{\omega^2}\le1.
\]
Therefore $0\le F_S(\lambda)\le2$ for every sample and every $\lambda>0$.
We state all constants used in the selection rule. Let $C_*\ge1$ be a universal constant for the smooth-loss theorem of Srebro, Sridharan, and Tewari~\cite[Theorem~1]{sst}. In the form used here, for a nonnegative $s_0$-smooth loss bounded by $b$ on $\mathcal F$, with probability at least $1-\eta$, simultaneously for $f\in\mathcal F$,
\begin{align}
 R(f)-\widehat R_S(f)\le C_*\bigg[&\sqrt{\widehat R_S(f)}
 \left(\sqrt{s_0}\log^{3/2}(2n)\Rad_n(\mathcal F)+\sqrt{\frac{b\log(1/\eta)}n}\right)\notag\\
 &+s_0\log^3(2n)\Rad_n(\mathcal F)^2+\frac{b\log(1/\eta)}n\bigg].
 \label{eq:sst}
\end{align}
The Rademacher complexity here is worst-case empirical complexity. Squared loss is $2$-smooth.

For $0<\delta<1$, put
\begin{align}
 M_2&=1+(2D_0+8)^2, & K_n&=2D_0\log2+4\log(6M_2n^4)+1,\notag\\
 c_0&=\frac18+\frac{D_0\log2}{4}, &
 J_n&=\left\lceil\log_2(4n+K_n+2)\right\rceil,\notag\\
 t_n&=\max\left\{1,\log\frac{4(J_n+1)}\delta\right\},\notag\\
 a_n&=\max\left\{\frac1{\sqrt n},\ 2C_*\left[
       \frac{\sqrt2K_H\log^{3/2}(2n)}{\sqrt n}+2\sqrt{\frac{t_n}n}\right]\right\}.
 \label{eq:stat-constants}
\end{align}
Set $D_0=1$ and recall $K_H$ from~Lemma~\ref{lem:capacity}. Select the unique positive solution of
\begin{equation}
 \widehat\lambda_S^2
   =256a_n^2\{2F_S(\widehat\lambda_S)+\widehat\lambda_S(K_n+1)\}.
 \label{eq:calibration}
\end{equation}
This one-dimensional rule uses the observed free energy, not teacher resources or the population risk. The learner predicts with the posterior at this single selected value.

\begin{lemma}[Existence, uniqueness, and measurability]
\label{lem:root}
The root in~\eqref{eq:calibration} exists, is unique, and is a measurable function of the sample. It satisfies
$\widehat\lambda_S\ge256a_n^2(K_n+1)\ge1/n$.
No upper bound on $\widehat\lambda_S$ is required.
\end{lemma}
\begin{proof}
The likelihood integral makes $F_S$ continuous, and $F_S(\lambda)/\lambda$ is decreasing. The map
\[
 \Psi_S(\lambda)=\frac{\lambda^2}{2F_S(\lambda)+\lambda(K_n+1)}
 =\frac{\lambda}{2F_S(\lambda)/\lambda+K_n+1}
\]
is therefore continuous and strictly increasing. It tends to zero at zero because it is at most $\lambda/(K_n+1)$, and tends to infinity because $F_S\le2$. It crosses $256a_n^2$ exactly once. The lower bound follows directly from the root equation. For each $t>0$, the event $\{\widehat\lambda_S<t\}$ equals $\{\Psi_S(t)>256a_n^2\}$, proving measurability.
\end{proof}

\begin{lemma}[One statistical event for all widths and scales]
\label{lem:stat-event}
With probability at least $1-\delta/2$, simultaneously for every supported network with $Q_N\le4n+K_n+1$,
\begin{equation}
 R(z_w)-\widehat R_S(z_w)
 \le a_n\sqrt{\widehat R_S(z_w)}(1+Q_N)+a_n^2(1+Q_N)^2.
 \label{eq:stat-event}
\end{equation}
\end{lemma}
\begin{proof}
Apply~\eqref{eq:sst} to $\mathcal F_{2^j}$, $0\le j\le J_n$, with failure probability $\delta/[2(J_n+1)]$. By~\eqref{eq:energy-lower}, its loss is bounded by $4(1+2^j)^2$, and by Lemma~\ref{lem:capacity}, its Rademacher complexity is at most $K_H(1+2^j)/\sqrt n$. Use the unit cutoff for $Q_N\le1$ and the least containing dyadic cutoff otherwise; then $1+2^j\le2(1+Q_N)$. The definition of $a_n$ and a union bound prove the claim. The energy-ball classes already include every width, so the event does not depend on $N$, $\gamma$, or $\lambda$.
\end{proof}

\begin{lemma}[Weighted empirical-objective tails]
\label{lem:tail}
For every sample and $\lambda>0$, the posterior raw-risk integral on
\[
 \{\widehat R_S+\lambda Q_N>2F_S(\lambda)+\lambda K_n\}
\]
is at most $n^{-4}$. Also, for every finite-entropy law $\rho$ relative to the prior,
\begin{equation}
 \int Q_Nd\rho\le\frac4\omega\KL(\rho\|\pi)+2D_0\log2.
 \label{eq:entropy}
\end{equation}
\end{lemma}
\begin{proof}
On $J:=\widehat R_S+\lambda Q_N>u$,
\[
 e^{-\omega\widehat R_S/(2\lambda)}
 \le e^{-\omega u/(4\lambda)}e^{\omega Q_N/4}.
\]
By~\eqref{eq:chi-square}, exponential tilting by $e^{\omega Q_N/4}$ doubles the chi-square scale and contributes the factor $2^{N(d+1)/2}$. Under that tilted law,
\[
 \E Q_N^2=\frac4{\omega^2}N(d+1)\{N(d+1)+2\}\le(2D_0+8)^2.
\]
After dividing by the posterior normalizer,
\[
 \int(1+Q_N^2)\mathbf1_{\{J>u\}}d\Pi_\lambda
 \le M_2\exp\left[\frac{\omega F_S}{2\lambda}-\frac{\omega u}{4\lambda}
                       +\frac{N(d+1)\log2}{2}\right].
\]
Raw risk is at most $2Q_N^2+2\le6(1+Q_N^2)$. Use $u=2F_S+\lambda K_n$, $N(d+1)\le\omega D_0$, and $\omega\ge1$ to obtain the claimed tail bound.

Finally $\log\int e^{\omega Q_N/4}d\pi=N(d+1)\log2/2\le\omega D_0\log2/2$. The entropy variational inequality gives~\eqref{eq:entropy}.
\end{proof}

\begin{theorem}[An observable certificate and an empirical oracle]
\label{thm:certificate}
On the event of Lemma~\ref{lem:stat-event}, simultaneously for every $N\ge1$ and $\gamma\ge1$,
\begin{equation}
 \int R(z_w)d\Pi_{\widehat\lambda_S}
 \le F_S(\widehat\lambda_S)+c_0\widehat\lambda_S+n^{-4}.
 \label{eq:certificate}
\end{equation}
If $F_S(\lambda)\le u+\lambda T$ for every positive $\lambda$, put
$\mathcal B=2T+K_n+1+c_0$. Then
\begin{equation}
 \int R(z_w)d\Pi_{\widehat\lambda_S}
 \le u+16\sqrt2\mathcal B a_n\sqrt u+256\mathcal B^2a_n^2+n^{-4}.
 \label{eq:empirical-oracle}
\end{equation}
Both bounds also hold for the ordinary raw posterior mean.
\end{theorem}
\begin{proof}
Write $\lambda=\widehat\lambda_S$ and $G=2F_S(\lambda)+\lambda(K_n+1)$. Off the tail in Lemma~\ref{lem:tail},
\[
 \widehat R_S\le G,\qquad1+Q_N\le G/\lambda,\qquad Q_N\le4n+K_n.
\]
The last inequality uses $F_S\le2$ and Lemma~\ref{lem:root}. The two terms of~\eqref{eq:stat-event} are bounded, respectively, by
\[
 a_n\sqrt G(1+Q_N)=\frac\lambda{16}(1+Q_N),
 \qquad
 a_n^2\frac G\lambda(1+Q_N)=\frac\lambda{256}(1+Q_N).
\]
Their sum is at most $\lambda(1+Q_N)/8$. Integrate and use~\eqref{eq:entropy}. The energy term costs at most one quarter of the relative-entropy penalty retained in $F_S$, plus $\lambda D_0\log2/4$. Adding $\lambda/8$ and the weighted tail gives~\eqref{eq:certificate}.

For the comparison, the root implies
\[
 \lambda^2\le512a_n^2u+256a_n^2(2T+K_n+1)\lambda.
\]
The elementary implication $x^2\le bx+c\Rightarrow x\le b+\sqrt c$ gives
$\lambda\le256a_n^2\mathcal B+16\sqrt2a_n\sqrt u$.
Substitute into~\eqref{eq:certificate} and use $T+c_0\le\mathcal B$ to obtain~\eqref{eq:empirical-oracle}. All raw moments exist by~\eqref{eq:energy-lower} and the Gaussian endpoint tails. Jensen gives the mean-prediction claim.
\end{proof}

\begin{theorem}[A dimension-free path-variation oracle]
\label{thm:main}
Fix any ReLU teacher $g$ with hidden depth $H$, width $C$, and mixed path variation $P$. Define $T_g$ by~\eqref{eq:teacher-budget}, and put $\mathcal B_g=2T_g+K_n+1+c_0$. With probability at least $1-\delta$, simultaneously for every $N\ge C$ and every $\gamma\ge1$,
\begin{align}
 R\left(\int z_wd\Pi_{\widehat\lambda_S}\right)
 \le\int R(z_w)d\Pi_{\widehat\lambda_S}
 \le R(g)&+25a_n\mathcal B_g\sqrt{R(g)}\notag\\
 &+280a_n^2\mathcal B_g^2+n^{-4}.
 \label{eq:population-oracle}
\end{align}
The calibration uses none of $C,P,R(g)$ and no coefficient bound is required. The data may have any IID law on $\mathcal X\times[-1,1]$.
\end{theorem}
\begin{proof}
The comparison~\eqref{eq:free-energy-comparison} supplies $u=\widehat R_S(g)+1/(2n)$. Put $E=4+2P$. Then $|g|\le P\le E$ and $T_g\ge E\ge4$. Bernstein gives
\[
 \widehat R_S(g)\le R(g)+(E+1)\sqrt{\frac{2R(g)\log(2/\delta)}n}
                  +\frac{2(E+1)^2\log(2/\delta)}{3n}
\]
except on an event of probability at most $\delta/2$. The inequalities $(E+1)^2\le2E^2$, $a_n^2\ge16\log(2/\delta)/n$, and $a_n^2\ge1/n$ imply
$u\le(\sqrt{R(g)}+a_nE)^2$.
Substitute this and its square root into~\eqref{eq:empirical-oracle}, using $E\le T_g\le\mathcal B_g$. The square-root coefficient is at most $2+16\sqrt2<25$, and the quadratic coefficient at most $257+16\sqrt2<280$. Combining the two sample events proves the result. Both are common to every eligible learner width and exponent.
\end{proof}

\begin{corollary}[Sample complexity in path variation]
\label{cor:polynomial}
For a comparison network with path variation $P$ and risk $R_*$, a sufficient sample count for oracle excess at most $\varepsilon\in(0,1)$ is
\begin{equation}
 n=\widetilde O\left(H(1+P)^2\left\{\frac{R_*}{\varepsilon^2}+\frac1\varepsilon\right\}\right).
 \label{eq:sample-budget}
\end{equation}
Only logarithms in $n,H,P,1/\varepsilon,1/\delta$ are suppressed. There is no input-dimension factor once $N\ge C$. The teacher resources and risk are used to state the guarantee, not to choose the posterior.
\end{corollary}
\begin{proof}
Here $a_n^2=\widetilde O(H/n)$ and $\mathcal B_g=\widetilde O(1+P)$. Substitute in~\eqref{eq:population-oracle}. These estimates follow directly from the displayed constants with $D_0=1$ and $K_H=1+\sqrt{2H\log2}$.
\end{proof}

\begin{corollary}[Quadratic width and linear depth under controlled gain]
\label{cor:width-depth}
Suppose a width-$C$ teacher has first-row Euclidean norms at most $B_1$, readout coefficients at most $B_2$, and
\[
 G=\prod_{\ell=2}^H\bigl\||W^\ell|\bigr\|_{\op}.
\]
Then $P\le GCB_1B_2$. In particular, at fixed $B_1,B_2$ and $G\le1$, the sufficient sample count is
\[
 \widetilde O\left(C^2H\left\{\frac{R_*}{\varepsilon^2}+\frac1\varepsilon\right\}\right).
\]
Only the product of gains is needed for this conclusion; the teacher need not have every gain at most one.
\end{corollary}
\begin{proof}
Apply Cauchy--Schwarz to~\eqref{eq:pathnorm}: $P\le\|v\|_2G\|a^1\|_2\le(\sqrt C B_2)G(\sqrt C B_1)$. Then use Corollary~\ref{cor:polynomial}.
\end{proof}

For arbitrary coefficient-bounded teachers, this estimate still gives a fixed-depth polynomial: if every internal coefficient and readout coefficient is at most $B$ in absolute value and first-row norms are at most $B_1$, then $P\le C^HB^HB_1$. Under a coordinatewise first-layer cap $B$ as well, $B_1\le\sqrt d B$, giving $P\le\sqrt d C^HB^{H+1}$. Its possible exponential dependence on depth is not removed by neuron rescaling.

\section{What this establishes and what it does not}
\paragraph{A function-prior statement.}
Theorem~\ref{thm:prior} bounds actual prior probabilities of function neighborhoods. Theorem~\ref{thm:energy-characterization} identifies the infinite-width minimum quadratic energy exactly with twice minimal path variation. These are stronger claims than calling a sample-complexity upper bound a prior. They still do not assert that every pair of functions is ranked exactly by path variation or that posterior computation is efficient.

\paragraph{The support condition changes parameter preferences.}
Absolute spectral nonexpansiveness is stronger than ordinary operator-norm nonexpansiveness: signed cancellations cannot hide amplification. Yet Lemma~\ref{lem:flow} shows that at a fixed width and depth it removes no ReLU functions. What changes is the parameterization and its Gaussian weight. General small networks may require large path variation, so this fact alone gives no unrestricted polynomial dependence jointly on teacher width and depth.

\paragraph{The Gaussian scaling is explicit.}
The first layer and readout have common precision $N(d+1)+(H-1)N^{\gamma+1}$, while internal coordinates have variance $N^{-\gamma}$. The likelihood is calibrated to the endpoint precision. This is a particular family of layerwise priors; the theorem does not transfer automatically to the more usual tuple with identical normalized strength at all layers. Its internal exponent can be any $\gamma\ge1$. More generally, internal precisions $\alpha_\ell>0$ may vary by layer: choosing $\omega=N(d+1)+N\sum_{\ell=2}^H\max\{N,\alpha_\ell\}$ and internal variances $1/\alpha_\ell$ gives the same results. Total dimension divided by $\omega$ is at most one, and the cloned internal mean cost is at most $C/N$; all other proofs are unchanged. This extension concerns calibrated endpoint priors, not an exhaustive claim about arbitrary unchanged scaling tuples.

\paragraph{Dimension and learner width.}
The statistical rate is dimension-free on the Euclidean unit ball, at every learner width above the comparison width. The architecture-dependent precision includes the first-layer parameter volume, so no additional lower bound on width in terms of input dimension is needed. The prior is specified separately at each width. If one starts with a cube and coordinatewise coefficient caps, normalizing the input and its constant coordinate changes first-row Euclidean norms; those factors must be counted in the teacher path variation. The theorem does not erase that distinction between norm and coordinate bounds.

\paragraph{Noise and approximation.}
Labels need not be generated by a network or by Gaussian noise. The likelihood is selected from the observed free energy and automatically gives a faster bound when the comparison risk is small. Prediction averages raw outputs. The result supplies neither an efficient free-energy calculation nor a gradient-descent guarantee.

\paragraph{Verification status.}
The arguments are conventional proofs using the published smooth-loss theorem with its assumptions stated explicitly. They have not been formalized in Lean.

\begin{thebibliography}{9}
\bibitem{sst} Nathan Srebro, Karthik Sridharan and Ambuj Tewari. \emph{Smoothness, Low Noise and Fast Rates}. Advances in Neural Information Processing Systems 23, 2010. \url{https://www.cs.cornell.edu/~sridharan/smoothness.pdf}.
\end{thebibliography}
\end{document}
